% chapters/ch09_theoretical_invariants.tex
% Chapter 9: Theoretical Invariants
% Source: NRMO_v5_updated.pdf, Chapter 9 (lines 6735-7090 of v5_layout.txt)

\chapter{Theoretical Invariants}
\label{ch:theoretical-invariants}

This chapter restates the v2.0 theoretical invariants in their final
form (UNCHANGED), establishes what remains theoretically guaranteed,
and documents the v2.1 operational refinements (NEW) that build on
them without weakening any invariant.

% =====================================================================
\section{Section 2: Theoretical Invariants (UNCHANGED)}
\label{sec:invariants-unchanged}

\subsection{Non-Ruin Constraint}
\label{sec:inv-non-ruin}

\begin{equation}
P(s_t \in R \mid \pi, \varepsilon) \le \varepsilon
\quad \text{for all } t,
\label{eq:non-ruin-constraint}
\end{equation}

\noindent
where:
\begin{itemize}
\item $R \subset S$: ruin states (absorbing, irreversible);
\item $\varepsilon$: externally supplied tolerance parameter;
\item $\pi$: policy (action selection mechanism);
\item $s_t$: state at time $t$.
\end{itemize}

\paragraph{Invariant status.}
This constraint is unchanged and remains the primary guarantee of
the framework.

\subsection{State Space}
\label{sec:inv-state-space}

\begin{equation}
S = \{\mathrm{ACTIVE}, \mathrm{HOLD}, \mathrm{EXIT}, \mathrm{SAFE\_EXIT}\}.
\end{equation}

\paragraph{Transitions.}
\begin{itemize}
\item $\mathrm{ACTIVE} \to \mathrm{HOLD}$: human intervention required.
\item $\mathrm{ACTIVE} \to \mathrm{EXIT}$: planned termination.
\item $\mathrm{ACTIVE} \to \mathrm{SAFE\_EXIT}$: emergency
termination.
\item $\mathrm{HOLD} \to \mathrm{ACTIVE}$: human approval.
\item $\mathrm{HOLD} \to \mathrm{SAFE\_EXIT}$: escalation.
\end{itemize}

\paragraph{Invariant status.}
State space and transitions unchanged. Note: v2.1 operationally
refines HOLD into subtypes
(Section~\ref{sec:hold-refinement}), but this is an
implementation-level distinction that does not alter the theoretical
state space.

\subsection{APCSO Output Structure}
\label{sec:inv-apcso}

\begin{itemize}
\item Output: unordered set of exactly three choices
$\{a_1, a_2, a_3\} \subset A$.
\item Constraints:
\begin{enumerate}
\item No ranking or preference ordering.
\item At least one option must be HOLD or EXIT.
\item All options satisfy
$P(\mathrm{Ruin} \mid a) \le \varepsilon$.
\end{enumerate}
\end{itemize}

\paragraph{Invariant status.}
Choice-set structure unchanged.

\subsection{Autonomy Preservation}
\label{sec:inv-autonomy}

Principles (unchanged from v2.0):

\begin{itemize}
\item No single recommendation (prevents authority collapse).
\item No ranking (preserves decision autonomy).
\item Human-in-the-loop required for selection.
\item System provides structure; human provides judgement.
\end{itemize}

\subsection{Time-Path Feasibility}
\label{sec:inv-time-path}

Non-ergodicity awareness (unchanged):

\begin{itemize}
\item Irreversible transitions are tracked.
\item Path-dependent constraints enforced.
\item Ruin states are absorbing (no recovery).
\end{itemize}

% =====================================================================
\section{Section 3: Operational Refinements (NEW)}
\label{sec:operational-refinements}

The following refinements are introduced in v2.1. Each is shown to be
\emph{theoretically compatible} with the invariants above.

\subsection{Parameter Renewal Protocol}
\label{sec:parameter-renewal}

\paragraph{Theoretical position.}
In the theoretical model, $\varepsilon$ is an externally supplied
parameter at each time step. The theory permits time-dependent
$\varepsilon(t)$ without modification.

\paragraph{Operational problem.}
In practice, $\varepsilon$ tends to ossify when no renewal process
exists. Initial conservative values become entrenched, preventing
adaptation to new circumstances.

\paragraph{Refinement: $\varepsilon$-validity period.}

\begin{lstlisting}[basicstyle=\ttfamily\small,frame=single]
epsilon_config = {
  'value': float,                  // Current epsilon value
  'ratified_date': timestamp,      // Last approval date
  'renewal_period': T_renewal,     // e.g., 20 years
  'expiry_date': ratified_date + T_renewal,
  'status': VALID | PENDING_RENEWAL | EXPIRED
}

# Renewal Protocol
if current_time > epsilon_config.expiry_date:
  if not renewal_completed:
    epsilon_config.status = EXPIRED
    system_state = HOLD
    require_renewal_process()
\end{lstlisting}

\paragraph{Renewal requirements.}

\begin{table}[ht]
\centering
\small
\begin{tabularx}{\linewidth}{lXl}
\toprule
\textbf{Phase} & \textbf{Required artefacts} & \textbf{Timeline} \\
\midrule
Phase 1 (Review)
& Ruin Avoidance Log (RAL); Opportunity Cost Log (OCL); current
$\varepsilon$ performance metrics
& $T_{\mathrm{expiry}} - 6$ months \\
Phase 2 (Deliberation)
& Analysis of RAL vs OCL; proposal for $\varepsilon$ adjustment;
multiple $\varepsilon$ candidates
& $T_{\mathrm{expiry}} - 3$ months \\
Phase 3 (Approval)
& Explicit human re-approval; documentation of rationale; new
$\varepsilon$ config instantiation
& Before $T_{\mathrm{expiry}}$ \\
\bottomrule
\end{tabularx}
\caption[$\varepsilon$ renewal protocol]{Three-phase renewal protocol
for $\varepsilon$ revision. The protocol is mandatory: an expired
$\varepsilon$ forces a HOLD state until renewal is completed.}
\label{tab:epsilon-renewal}
\end{table}

\paragraph{Theoretical compatibility.}
The renewal protocol is theoretically compatible because
$\varepsilon(t)$ as a time-varying parameter is already permitted by
the theoretical model. The renewal protocol is a governance
mechanism for $\varepsilon$ supply; the non-ruin constraint
$P(\mathrm{Ruin} \mid \pi, \varepsilon) \le \varepsilon$ remains
enforced at all times. During EXPIRED status, the system defaults to
HOLD (safe fallback).

\subsection{HOLD State Refinement}
\label{sec:hold-refinement}

\paragraph{Theoretical position.}
HOLD is defined as a state requiring human intervention. The
theoretical model does not specify exit conditions or internal
structure.

\paragraph{Operational problem.}
Without operationalised exit conditions, HOLD states can persist
indefinitely, becoming a de facto rejection mechanism rather than
a temporary pause.

\paragraph{Refinement: HOLD subtypes.}

\begin{lstlisting}[basicstyle=\ttfamily\small,frame=single]
# Theoretical level (unchanged)
State = {ACTIVE, HOLD, EXIT, SAFE_EXIT}

# Operational level (refined)
HOLD -> {
  HOLD_WITH_OBSERVATION,    # Sensor enhancement mode
  HOLD_WITH_SHADOW,         # Delegate to simulation
  HOLD_WITH_EXPIRY          # Explicit timeout
}

# All HOLD subtypes enforce finite duration
\end{lstlisting}

\begin{table}[ht]
\centering
\small
\begin{tabularx}{\linewidth}{lXXl}
\toprule
\textbf{Subtype} & \textbf{Purpose} & \textbf{Exit condition} & \textbf{Max duration} \\
\midrule
HOLD\_WITH\_OBSERVATION
& Information gathering
& Sensor data collected
& 30 days \\
HOLD\_WITH\_SHADOW
& Risk assessment
& Shadow simulation complete $\to$ REVIEW
& Shadow duration \\
HOLD\_WITH\_EXPIRY
& External condition wait
& Timeout $\to$ AUTO\_EXIT or HUMAN\_DECISION
& 90 days \\
\bottomrule
\end{tabularx}
\caption[HOLD subtypes]{HOLD subtype definitions with exit conditions
and maximum durations. The maximum-duration constraint prevents
indefinite HOLD.}
\label{tab:hold-subtypes}
\end{table}

\paragraph{Selection rules.}

\begin{lstlisting}[language=Python,basicstyle=\ttfamily\small,frame=single]
def select_hold_type(reason):
    if information_gap_identified and observable_metric_exists:
        return HOLD_WITH_OBSERVATION
    elif risk_assessment_needed and simulation_feasible:
        return HOLD_WITH_SHADOW
    elif external_condition_pending:
        return HOLD_WITH_EXPIRY
    else:
        # Indefinite HOLD not permitted
        escalate_to_human_decision()
\end{lstlisting}

\paragraph{Prohibited patterns.}
\begin{itemize}
\item Unconditional HOLD (no exit condition specified).
\item Infinite HOLD (no timeout).
\item HOLD chaining (HOLD $\to$ HOLD without resolution).
\item HOLD abuse (same issue HOLD'd 3+ times).
\end{itemize}

\paragraph{Theoretical compatibility.}
HOLD remains a single abstract state in the theoretical model.
Subtypes are implementation-level distinctions. The state-transition
graph $S = \{\mathrm{ACTIVE}, \mathrm{HOLD}, \mathrm{EXIT},
\mathrm{SAFE\_EXIT}\}$ is unchanged. Exit-condition enforcement is a
deadlock-prevention mechanism (implementation concern), not a
theoretical relaxation.

\subsection{Zero-Risk Fast Path}
\label{sec:zero-risk-fast-path}

\paragraph{Theoretical position.}
Actions satisfying $\mathrm{irreversibility} = 0$ and
$\varepsilon\text{-consumption} = 0$ have $P(\mathrm{Ruin}) = 0$ by
definition. Formal verification is theoretically unnecessary for
such actions.

\paragraph{Operational opportunity.}
Zero-risk actions (e.g., art projects, local experiments, reversible
cultural activities) can be routed through a computationally
optimised path, as the constraint check is trivial.

\paragraph{Refinement: fast-path routing.}

\begin{lstlisting}[language=Python,basicstyle=\ttfamily\small,frame=single]
def process_action(action):
    # Fast path condition check
    if (action.irreversibility == 0 and
        action.epsilon_consumption == 0 and
        action.scope == "local"):
        # P(Ruin) = 0 trivially satisfied
        # Skip formal verification (optimization)
        route_to_human_decision(action, fast_path=True)
    else:
        # Standard NRMO verification path
        verify_constraints(action)
        if passed:
            route_to_human_decision(action, fast_path=False)
\end{lstlisting}

\paragraph{Eligibility criteria.}

\begin{table}[ht]
\centering
\small
\begin{tabularx}{\linewidth}{lXX}
\toprule
\textbf{Criterion} & \textbf{Requirement} & \textbf{Example (eligible / ineligible)} \\
\midrule
Irreversibility
& $= 0$ (fully reversible)
& Theatre performance / building construction \\
$\varepsilon$-consumption
& $= 0$ (zero ruin risk)
& Poetry reading / financial product launch \\
Scope
& Local \& contained
& Community festival / national policy change \\
Termination
& Explicit endpoint
& 3-month experiment / permanent modification \\
\bottomrule
\end{tabularx}
\caption[Fast-path eligibility]{Fast-path eligibility criteria. All
four must hold; otherwise the action is routed through standard
NRMO verification.}
\label{tab:fast-path-eligibility}
\end{table}

\paragraph{Distinction from Shadow Optimization.}

\begin{table}[ht]
\centering
\small
\begin{tabularx}{\linewidth}{lXXX}
\toprule
\textbf{Mechanism} & \textbf{Risk level} & \textbf{Execution} & \textbf{Purpose} \\
\midrule
Shadow
& $\varepsilon > 0$ possible
& Simulated (isolated)
& High-risk experimentation \\
Fast Path
& $\varepsilon = 0$ (guaranteed)
& Real (fast-tracked)
& Zero-risk creative action \\
\bottomrule
\end{tabularx}
\caption[Fast Path vs Shadow]{Fast Path is distinct from Shadow
Optimisation. Shadow handles non-zero risk via simulation; Fast Path
handles zero risk by trivially passing the constraint check.}
\label{tab:fast-path-vs-shadow}
\end{table}

\paragraph{Theoretical compatibility.}
$P(\mathrm{Ruin}) = 0$ for $\varepsilon = 0$ and irreversibility
$= 0$ is mathematically trivial. The fast path is a computational
optimisation, not a constraint relaxation. The non-ruin constraint
$P(\mathrm{Ruin}) \le \varepsilon$ is satisfied by $0 \le 0$.
Fast-path actions still require human approval (autonomy preserved).

\subsection{Abuse Pattern Detection}
\label{sec:abuse-detection}

\paragraph{Operational problem.}
NRMO terminology can be misused for non-ruin-related blocking,
e.g., ``This is NRMO-incompatible'' used to halt legitimate debates.

\paragraph{Refinement: heuristic abuse-pattern detection.}

\begin{lstlisting}[basicstyle=\ttfamily\scriptsize,frame=single]
abuse_patterns = {
  'perpetual_hold': {
    'detection': hold_duration > 90_days
                  and exit_condition == null,
    'severity': CRITICAL,
    'action':   escalate_to_shadow_review
  },
  'safe_mode_normalization': {
    'detection': safe_mode_ratio > 0.8 for 180_days,
    'severity': HIGH,
    'action':   trigger_epsilon_renewal
  },
  'vocabulary_misuse': {
    'detection': nrmo_invoked_for_non_ruin_argument,
    'severity': MEDIUM,
    'action':   log_and_flag_for_review
  },
  'risk_free_rejection': {
    'detection': (action.rejected
                  and action.epsilon_consumption == 0
                  and action.reversible == true),
    'severity': HIGH,
    'action':   require_explicit_justification
  }
}
\end{lstlisting}

\paragraph{Detection and response.}

\begin{table}[ht]
\centering
\small
\begin{tabularx}{\linewidth}{lXX}
\toprule
\textbf{Pattern} & \textbf{Automated action} & \textbf{Human notification} \\
\midrule
Perpetual HOLD
& Force Shadow Review
& Immediate \\
SAFE\_MODE Normalisation
& Trigger $\varepsilon$ renewal process
& Within 7 days \\
Vocabulary Misuse
& Log \& monitor
& Monthly report \\
Risk-Free Rejection
& Request justification
& Immediate \\
\bottomrule
\end{tabularx}
\caption[Abuse patterns detection and response]{Abuse-pattern
heuristic detection rules and corresponding automated responses.}
\label{tab:abuse-patterns}
\end{table}

\paragraph{Theoretical compatibility (limitations acknowledged).}
\begin{itemize}
\item Abuse detection relies on heuristics, not formal guarantees.
\item Not part of the theoretical model (purely operational
safeguard).
\item Requires human judgement for final escalation decisions.
\item False positives possible; operates as monitoring, not
enforcement.
\end{itemize}

% =====================================================================
\section{Section 4: Deployment Recommendations}
\label{sec:deployment-recommendations}

\subsection{Priority Implementation Order}

\begin{table}[ht]
\centering
\small
\begin{tabularx}{\linewidth}{lXXl}
\toprule
\textbf{Priority} & \textbf{Refinement} & \textbf{Rationale} & \textbf{Difficulty} \\
\midrule
Critical
& HOLD State Refinement
& Prevents indefinite stasis
& Medium \\
Critical
& Abuse Pattern Detection
& Prevents internal misuse
& Medium \\
Important
& Parameter Renewal Protocol
& Prevents $\varepsilon$ ossification (long-term)
& Low \\
Recommended
& Zero-Risk Fast Path
& Enables low-risk creativity
& High \\
\bottomrule
\end{tabularx}
\caption[Deployment priority]{Recommended deployment priority for
the v2.1 operational refinements.}
\label{tab:deployment-priority}
\end{table}

\subsection{Integration with Existing Systems}

\begin{lstlisting}[language=Python,basicstyle=\ttfamily\small,frame=single]
# v2.0 -> v2.1 Migration Path

# 1. Extend epsilon structure
epsilon_v3 = {
    **epsilon_v2,
    'renewal_period': T_renewal,
    'expiry_date':    compute_expiry(),
    'ral': [],   # Ruin Avoidance Log
    'ocl': []    # Opportunity Cost Log
}

# 2. Refine HOLD state
if state == HOLD:
    state = select_hold_subtype(context)
    enforce_exit_condition(state)

# 3. Add fast path router
if is_zero_risk(action):
    fast_path_route(action)
else:
    standard_nrmo_verification(action)

# 4. Enable abuse monitoring
abuse_monitor.start()
\end{lstlisting}

\subsection{Institutional Requirements}

\paragraph{Required institutional capabilities.}

\begin{itemize}
\item \textbf{$\varepsilon$ Governance Body}: authority to approve
$\varepsilon$ renewals.
\item \textbf{Shadow Infrastructure}: for HOLD\_WITH\_SHADOW
escalations.
\item \textbf{Monitoring Team}: to review abuse-pattern alerts.
\item \textbf{Audit Trail}: comprehensive logging of all decisions.
\end{itemize}

\subsection{Backward Compatibility}

v2.1 is backward compatible with v2.0:

\begin{itemize}
\item All v2.0 theoretical guarantees preserved.
\item Existing deployments can adopt refinements incrementally.
\item No breaking changes to core interfaces.
\item v2.0 behaviour recoverable by disabling operational
refinements.
\end{itemize}

% =====================================================================
\section{Section 5: Limitations and Future Work
(ACKNOWLEDGED)}
\label{sec:invariants-limitations}

\subsection{$\varepsilon$ Supply Mechanism}

\paragraph{Current state.}
The theoretical model treats $\varepsilon$ as externally supplied.
The renewal protocol specifies \emph{when} to update $\varepsilon$
but not \emph{how} consensus is reached.

\paragraph{Open questions.}
\begin{itemize}
\item Who has authority to approve $\varepsilon$ changes?
\item How are conflicting stakeholder preferences resolved?
\item What mechanisms ensure $\varepsilon$ reflects legitimate risk
tolerance?
\end{itemize}

\paragraph{Future work.}
Formalise democratic $\varepsilon$-governance protocols, potentially
integrating with existing institutional decision-making frameworks.

\subsection{HOLD Exit Conditions}

\paragraph{Current state.}
HOLD subtypes enforce finite duration, but ultimately rely on human
decision to exit. The theoretical model does not specify what
happens if humans never decide.

\paragraph{Mitigation.}
Operational timeouts force escalation, but this is an implementation
choice, not a theoretical guarantee.

\paragraph{Future work.}
Explore formal models of human-in-the-loop decision latency and
develop escalation protocols for pathological cases.

% =====================================================================
\section*{Connections to Other Parts}

\begin{itemize}
\item Equation~\ref{eq:non-ruin-constraint} (non-ruin constraint) is
the v2.1 form of the survival-first axiom NRMO-0
(Section~\ref{sec:foundations-3-1-definition}).
\item APCSO output structure
(Section~\ref{sec:inv-apcso}) is the formal core of the
choice-set system in
Section~\ref{sec:apcso-formal}.
\item HOLD subtypes
(Section~\ref{sec:hold-refinement}) correspond to the
partial-observability handling in the simulator (Part~VIII):
HOLD\_WITH\_OBSERVATION is the cognitive analogue of MAPLayer's
NEUTRAL/PROBE modes.
\item The $\varepsilon$ renewal protocol
(Section~\ref{sec:parameter-renewal}) is the cognitive analogue of
the Adaptive Tuning Layer's hysteresis mechanism (Part~VIII).
\end{itemize}
